\section{Where the budget goes}\label{sec:budget}
\noindent LLM-NAS states its budget in LLM calls: 120 per search, of which 60 ask for four children each, so 240 architectures are requested. The framework does not replace a rejected proposal, so the number of architectures actually evaluated is 30 initial ones plus the accepted children, and it differs from method to method. This section measures how many proposals are lost, why, and where the lost ones come from. All numbers are for CIFAR-10 on six devices (one seed per LLM; Claude Opus~5.5 on Edge GPU only; three seeds for the controls). They were obtained by replaying the archived prompts and replies (\code{python -m llmnas duplicates}) and match the counts recorded during the searches exactly; no LLM was called.

\paragraph{How proposals are rejected.} A proposal is checked in a fixed order and rejected at the first failure: the reply cannot be parsed; the cell is not a legal NAS-Bench-201 code (invalid); it breaks the niche's operator counts (out of niche); it has already been evaluated in this search (duplicate); or HW-NAS-Bench has no latency for it on the device (in practice, the Edge TPU cells that failed to compile). Table~\ref{tab:waste} gives the outcome of every proposal. No reply failed to parse. Duplicates are the largest loss for every LLM except Qwen2.5-7B, which loses about as many proposals to broken niche constraints, as does DeepSeek-Coder-V2. The result is that the LLM runs evaluate 76--222 architectures, against 270 for NSGA-II \cite{deb02} and regularised evolution \cite{real19} run with the same evaluator, because those two algorithms keep generating until the budget is filled, as is standard in evolutionary computation \cite{mauldin84}. Random-operator evolution (PEA) inside the same framework is not refilled either and evaluates 169 on average. Its duplicate rate (40\%) is close to GPT-4.1's (43\%), so for the paper's own model much of the loss is a property of the framework, which does not refill, rather than of the language model; the weaker models lose more.

\begin{table}[h]\centering\footnotesize
\caption{Outcome of every proposal, CIFAR-10, six devices (shares of all proposals), and the mean number of architectures evaluated per search, $N$ (30 initial plus accepted children). The right-hand block splits the duplicates by where the repeated cell had been seen: one of the parents listed in the same prompt; one of the previous cycle's children listed in the same prompt; an earlier child in the same reply; or a cell evaluated in an earlier generation and no longer shown. Controls: three seeds. NSGA-II and regularised evolution regenerate until $N=270$.}\label{tab:waste}
\setlength{\tabcolsep}{3.5pt}
\resizebox{\linewidth}{!}{\begin{tabular}{l|cccc|c|cccc}
\toprule
 & \multicolumn{4}{c|}{Outcome of proposals (\%)} & & \multicolumn{4}{c}{Duplicates by source (\% of proposals)} \\
Method & new & duplicate & out of niche & other & $N$ & parent shown & last cycle shown & same reply & earlier, not shown \\
\midrule
GPT-4.1                     & 47.9 & 42.7 & 8.3  & 1.1 & 145 & 5.8  & 2.8  & 1.7 & 32.3 \\
Claude Haiku~4.5            & 23.3 & 65.6 & 9.3  & 1.8 & 86  & 32.5 & 3.8  & 0.6 & 28.6 \\
Claude Opus~5.5 (Edge GPU)  & 80.0 & 20.0 & 0.0  & 0.0 & 222 & 0.0  & 0.0  & 0.0 & 20.0 \\
Qwen2.5-7B (local)          & 38.8 & 28.1 & 32.1 & 1.1 & 123 & 11.7 & 6.1  & 0.5 & 9.7 \\
DeepSeek-Coder-V2-16B (local) & 19.0 & 57.1 & 22.4 & 1.5 & 76 & 41.1 & 12.6 & 0.3 & 3.1 \\
\midrule
PEA (random operators)      & 58.0 & 39.8 & --   & 2.2 & 169 & \multicolumn{4}{c}{--} \\
Random sampling in niches   & 99.7 & 0.3  & --   & 0.0 & 269 & \multicolumn{4}{c}{--} \\
NSGA-II                     & 74.7 & 17.6 & --   & 7.7 & 270 & \multicolumn{4}{c}{--} \\
Regularised evolution       & 76.8 & 20.5 & --   & 2.7 & 270 & \multicolumn{4}{c}{--} \\
\bottomrule
\end{tabular}}

\smallskip
\raggedright\scriptsize ``Other'' = no latency on the device, plus invalid codes (at most 0.9\% for the LLMs; 4.6\% for NSGA-II, whose crossover of two near-identical parents can fail to produce a new child). The controls draw children inside the niche by construction (``--'').
\end{table}

\paragraph{Where the duplicates come from.} The models fail in two different ways. The weaker models copy what they were shown: 41\% of DeepSeek-Coder-V2's proposals and 33\% of Haiku's are one of the parents in the same prompt, returned unchanged although a mutation or crossover was requested. DeepSeek does this from the first generation (69\% duplicates in generation~1). GPT-4.1 rarely copies a parent (6\%). Its duplicates are cells evaluated in an earlier generation that are no longer in the prompt (32\% of its proposals). Every one of these cells had appeared in an earlier prompt, a median of three to five generations before, but each call is stateless and the prompt shows only six parents and the last cycle's children, so the model cannot see that the design has already been tried and derives it again. The duplicate rate therefore grows as the search proceeds: for GPT-4.1 it rises from 3\% in generation~1 to 65\% in generations 8--10, for Haiku from 5\% to over 80\%, and even Claude Opus~5.5 reaches 42--58\% in the last two generations. Running out of space is not the cause. When a proposal was made, on average only 0.2--4.5\% of its niche had been evaluated; even the smallest niche, S5 (265 cells), can be at most 17\% explored by the end of a search (5 initial cells and 40 children). The duplicate rate is 39--53\% in every niche, and it is not higher in the small niches (S5: 40\%; S2, 6{,}144 cells: 52\%). Crossover produces more duplicates than mutation for every model except Opus (GPT-4.1: 47\% against 40\%), because two similar parents often recombine into one of themselves. The pattern is the narrowing that the paper calls exploration bias, and it matches what trajectory studies of LLM-guided evolution report: strong models act as local refiners and narrow the search semantically \cite{zhang26traj}, and LLM optimisers converge to the same solutions regardless of the problem instance or the temperature \cite{redko26}.

\paragraph{What the loss costs in front quality.} Because the LLM runs evaluate fewer architectures, part of their lower hypervolume is a smaller sample rather than worse choices. The two parts can be separated by reading every method's trajectory at the same number of evaluations (Table~\ref{tab:matched}). After 145 evaluations, GPT-4.1 (0.928) is level with random sampling (0.930) but below PEA (0.944), NSGA-II (0.949) and regularised evolution (0.960) (means over devices, controls read at 145 evaluations). Between 145 and 270 evaluations random sampling and NSGA-II gain 0.020 and regularised evolution 0.007. If GPT-4.1 gained as much with a full budget it would be level with PEA and random sampling, and still below the two evolutionary algorithms. Our reading is that the smaller budget explains most of GPT-4.1's gap to random sampling, but not its gap to conventional evolution, which is a gap in the quality of its proposals. This is an extrapolation, to be tested by the refilling experiment in Section~\ref{sec:future}.

\paragraph{Relation to other work.} Filtering repeated candidates is common in LLM-driven search, but the studies we found report it differently. Work that generates network code removes near-duplicates by MinHash--Jaccard similarity and reports how many architectures were admitted \cite{khalid26,adhikari26}. A spiking-network search removes 30\% of raw LLM proposals as near-duplicates by a behavioural fingerprint \cite{zhao26spiking}. Program-evolution systems reject candidates by embedding similarity \cite{shinka}, or discard a duplicate and ask again up to a fixed number of attempts \cite{lyria}. None of these break duplicates down by where the repeated design came from, and they do not measure how the loss shrinks the evaluation budget relative to a baseline. The two LLM-based methods that report NAS-Bench-201 results do not analyse duplicates: GENIUS feeds every result back to GPT-4 and compares against random sampling after ten iterations \cite{genius}, and LLMatic generates 100 offspring per generation without a stated deduplication step \cite{llmatic}. The closest budget-matched comparison we found is in hyperparameter optimisation: an LLM advisor was overtaken by random search and Bayesian optimisation within a few evaluations, and its early lead came from a default configuration rather than from the LLM \cite{rodrigues26}. To our knowledge, this is the first exact count of duplicate LLM proposals in an enumerable NAS space, broken down by source, together with its effect on the evaluated budget. The remedies are known: rejecting and regenerating duplicates as evolutionary algorithms do \cite{mauldin84,lyria}, or showing the model its search history as OPRO does \cite{opro}. We report them as future work (Section~\ref{sec:future}), not as contributions.

\paragraph{Recommendation.} A comparison between an LLM operator and any other operator should be stated in evaluated architectures, not in LLM calls, and should report the share of proposals that were rejected. Under a budget of 120 calls, the paper's own model evaluates about half as many architectures as a budget-matched evolutionary algorithm.
